\begin{table*}[t]
\centering
\caption{Clean-channel reconstruction on 300 LibriSpeech test-clean utterances. Entries are means [95\% utterance-bootstrap intervals]. The first three codecs share the nominal 1.5-kbit/s operating point; bold means identify the best in that group. Opus and PCM are higher-rate references. Lower mel-L1 and higher STOI/PESQ-WB are better.}
\label{tab:rv2-quality}
\small\setlength{\tabcolsep}{5pt}
\begin{tabular}{lrrrr}
\toprule
Codec & kbit/s & mel-L1 & STOI & PESQ-WB\\
\midrule
SCIT & 1.5 & \textbf{0.864} [0.855, 0.874] & \textbf{0.892} [0.889, 0.895] & \textbf{2.097} [2.067, 2.128]\\
DAC ($n_q=3$) & 1.5 & 1.232 [1.219, 1.245] & 0.799 [0.796, 0.803] & 1.284 [1.266, 1.303]\\
EnCodec & 1.5 & 1.282 [1.266, 1.297] & 0.848 [0.845, 0.851] & 1.612 [1.585, 1.638]\\
Opus & 6.0 & 2.302$^\dagger$ [2.275, 2.330] & 0.906 [0.903, 0.908] & 2.237 [2.187, 2.287]\\
PCM & 256 & 0.000 [0.000, 0.000] & 1.000 [1.000, 1.000] & 4.644 [4.644, 4.644]\\
\bottomrule
\end{tabular}
\par\smallskip
\parbox{\textwidth}{\footnotesize $^\dagger$Full-band mel-L1 is sensitive to low-rate Opus bandwidth reduction and is not a standalone cross-codec perceptual ranking. Intervals use 10,000 resamples and seed 42; displayed values are rounded to three decimals.}
\end{table*}
